\section{Discussie}
\subsection{Aard en begrenzing van de bijdrage}
De voorgestelde bijdrage is een nog te toetsen contextbindingsprofiel. De brede integratie van masterdata en semantische technologie is reeds bestaand werk. \citep{L18} Het feit dat het profiel meer onderwerpen omvat, bewijst niet dat het een nieuw of beter ontwerp is. De afzonderlijke bouwstenen zijn grotendeels bestaand: terminologie, metadata\-registratie, provenance, catalogisering en ontologische analyse hebben elk eigen literatuur en specificaties. \citep{S11,S08,S33,S47,L08} De paper voegt een expliciete typering van de verbindingen, een releasegebonden afhankelijkheidsmodel en een toetsbare scheiding van gezag en afleiding toe als samenhangend voorstel. De wetenschappelijke waarde van die combinatie hangt af van vergelijking met verwante ontwerpen en toekomstige evaluatie; nieuwheid wordt niet afgeleid uit het aantal verbonden standaarden.

De reference architecture is bewust logischer dan fysiek. Een technische implementatie kan RDF centraal opslaan of semantische beschrijvingen federatief verbinden. Beide kunnen aan hetzelfde principe voldoen of op dezelfde requirement falen. Deze scheiding helpt voorkomen dat een reeds gebouwd prototype de literatuursynthese bepaalt. Zij sluit aan bij het onderscheid tussen artefact en overdraagbare ontwerpkennis in DSR. \citep{L04}

\subsection{Semantische volledigheid en de grenzen van formalisering}
Volledige explicitering van alle betekenis is geen realistische toetsdefinitie binnen dit onderzoek. Het ontwerp kiest daarom voor taakgebonden volledigheid: voldoende expliciete context om afgebakende vragen correct te kunnen beantwoorden en grenzen zichtbaar te houden. Deze keuze is een operationalisering, geen empirische vaststelling. De afstand tussen juridische tekst en formele regel blijft mogelijk bestaan; die afstand moet worden beschreven in plaats van door een succesvolle validator te worden verhuld.

Ook formele consistentie is beperkt. Een ontologie kan onder haar axioma’s consistent zijn terwijl de gekozen categorieën ongeschikt zijn voor de casus. Een dataset kan aan shapes voldoen terwijl de bronwaarde feitelijk onjuist is. De onderscheiden scopes van OWL, SHACL en kwaliteitsmodellen ondersteunen daarom afzonderlijke toetslagen. \citep{S30,S32,S12} Het SMDP beoogt deze beperkingen zichtbaarder te maken, niet op te heffen door één techniek.

\subsection{Gezag en federatie}
Een verbindingsgraaf kan institutionele verantwoordelijkheden representeren, maar niet creëren. Bij conflicterende definities van verschillende bevoegde bronnen moet de architectuur context en verschil kunnen bewaren. ``Eén waarheid'' is daarom geen ontwerpcriterium; wel wordt gevraagd of een consument de juiste uitspraak voor een afgebakende context en tijd kan identificeren. De basisregistratie-eisen en het FDS-kader motiveren aandacht voor verantwoordelijkheden, gebruik en terugmelding, maar vormen geen bewijs dat alle begrippen centraal moeten worden geharmoniseerd. \citep{S25,S27}

Hieruit volgt een open ontwerpafweging. Centrale standaardisatie kan uniformiteit in beschrijvingen nastreven, terwijl federatief beheer domeinspecifieke verantwoordelijkheid bewaart. Het voorgestelde manifest en de getypeerde mapping moeten beide kunnen ondersteunen. De kosten van coördinatie en onderhoud zijn nog onbekend en behoren expliciet tot de evaluatie.

\subsection{AI: theoretische motivatie zonder effectoverschatting}
De redenering dat meer automatische interpretatie expliciete semantiek belangrijker maakt, is in deze paper een conditionele ontwerpstelling. FAIR ondersteunt machinegericht hergebruik, maar een AI-systeem moet de aangeboden context nog correct selecteren en toepassen. \citep{L11} Kennisgrafen en retrievaltechnieken maken onderdelen van dat proces mogelijk, maar leveren geen algemene garantie op juiste of juridisch toelaatbare antwoorden. \citep{L13,L16}

R23 en DP9 richten zich daarom niet alleen op méér context, maar ook op het zichtbaar maken van onzekerheid, afleiding en begrensd gezag. Een foutieve machineleesbare annotatie kan juist systematisch worden verspreid. De toekomstige vergelijking moet zowel verbeteringen als nieuwe foutpatronen meten. AI blijft daarmee een veeleisende consument van hetzelfde dataproduct, geen rechtvaardiging om de overige gebruikers of governance te reduceren tot een AI-interface.

\section{Threats to Validity}
\textbf{Constructvaliditeit.} Begrippen als semantische volledigheid, ontologische consistentie en uitlegbaarheid kunnen verschillend worden geoperationaliseerd. Het voorstel beperkt ze tot afgebakende taken en maakt deelmaten zichtbaar. Deskundigenreview en negatieve gevallen moeten toetsen of de maten het bedoelde construct raken. Een volledig ingevuld metadataveld is niet automatisch inhoudelijk juist.

\textbf{Bron- en selectiebias.} Het corpus is doelgericht en deels gebaseerd op abstracts, bibliografische records en normcatalogi. Relevante integratieontwerpen kunnen ontbreken. De kennislacune blijft daarom voorlopig. Een systematische uitbreiding en toegang tot volledige normteksten zijn nodig voordat exclusieve nieuwheid of clausuleconformiteit wordt geclaimd. Meervoudige vermeldingen van dezelfde bron leveren geen onafhankelijke bevestiging.

\textbf{Interne validiteit.} In de A/B-vergelijking kunnen extra informatie, interfacekwaliteit, latency, retrievalkwaliteit en domeinervaring effecten verklaren. Gelijke broninhoud, vergelijkbare toegang, randomisering en gerichte ablaties moeten deze factoren zoveel mogelijk onderscheiden. Het uitvoerbare protocol moet nog aantonen dat deze controle haalbaar is.

\textbf{Onderzoekers- en prototypebias.} Ontwerpers kunnen hun eigen representatie gunstiger beoordelen of requirements achteraf op bestaande functies aanpassen. Vooraf vastgelegde criteria, onafhankelijke beoordeling, een gouden standaard buiten de implementatie en versiebevriezing moeten dit risico beperken. Het prototype wordt in deze paper niet als bewijs voor de architectuur gebruikt.

\textbf{Externe validiteit.} Het onderbouwde BAG-fragment representeert niet alle basisregistraties of overheidsdomeinen. MDM-literatuur uit bedrijven en retrievalonderzoek uit NLP-benchmarks kunnen niet rechtstreeks worden gegeneraliseerd naar publieke registratieketens. \citep{L05,L17,L16} Meerregistratiecases en praktijkgerichte evaluaties zijn noodzakelijk.

\textbf{Temporele validiteit.} Normen, profielen, lijststatussen en wetgeving hebben afzonderlijke versies en kunnen veranderen. Het corpus is op een peildatum afgebakend. De juridische en standaardtoepasselijkheid moet vóór uitvoering opnieuw worden vastgesteld; deze paper voorspelt die toekomstige status niet.

\textbf{Conclusievaliditeit.} Kleine taaksets, afhankelijke observaties en veel metrieken kunnen schijnbare effecten opleveren. De evaluatie moet primaire uitkomsten vooraf vastleggen, afhankelijkheden respecteren en onzekerheid rapporteren. Er zijn nog geen gegevens voor een effectgrootte of steekproefberekening.

\textbf{Ecologische validiteit en beheerlast.} Synthetische foutgevallen kunnen eenvoudiger zijn dan langdurig beheer met onvolledige bronnen, conflicterende belangen en distributiestoringen. De voorgestelde praktijkepisode en kostenmeting zijn daarom noodzakelijk. Formele correctheid in een laboratorium geldt niet als bewijs van organisatorische uitvoerbaarheid.

\textbf{Bewijsvertrouwen.} Provenance en gezagsmetadata kunnen zelf onjuist zijn. Het model maakt claims herleidbaar, maar een bevoegde beoordeling van bronnen blijft nodig. Een persistent identifier of cryptografisch herkenbare afzender mag in de analyse niet worden verwisseld met ware inhoud. \citep{S33,S21,S39}

\textbf{Institutionele pluraliteit en toegangsbias.} Een centraal semantisch besluit kan legitieme perspectiefverschillen verbergen; een federatieve implementatie kan afhankelijkheidsfalen introduceren. Een gouden standaard kan de voorkeur van één bronhouder institutionaliseren. Beoordelaars moeten daarom ook meerdere toelaatbare antwoorden, conflicterende gronden en ontoegankelijke context kunnen registreren. De identiteit van de bron is geen automatische oplossing voor normatieve onenigheid.

\textbf{Interventielekkage en modellering als training.} Bij het bouwen van B kan extra domeinkennis ontstaan die beoordelaars of gebruikers vervolgens ook in A benutten. Een ontologievariant kan het antwoord verraden via haar labels. De voorbereidingstijd, expertise en informatie worden geregistreerd en de summatieve taakset blijft buiten het ontwerpteam. Een representatie die alleen door de eigen auteurs correct kan worden gebruikt, ondersteunt geen brede interoperabiliteitsclaim.

\textbf{Afhankelijkheid van de review en onvolledige bronverificatie.} De kritische review is door dezelfde AI-assistent uitgevoerd als de oorspronkelijke tekst. Zij is een gedocumenteerde zelfcorrectie, geen externe bevestiging. De audit markeert abstracts, catalogi, redirects en niet opnieuw inhoudelijk geverifieerde bronnen apart. Een succesvol opgehaalde pagina wordt niet als gelezen volledige norm gerekend.

\section{Conclusie}
Deze paper stelt een Semantic Master Data Product voor als een beheerde verbinding tussen betekenis, gezag, identiteit, historie, masterdata en distributie. De geselecteerde literatuur en standaarden bieden bouwstenen op verschillende abstractieniveaus. De ontwerpuitdaging ligt vooral in de relaties ertussen: begrip is geen klasse, data-element is geen waarde, record is geen entiteit, provenance is geen authenticiteit en catalogusmetadata vormen niet zelfstandig een volledig productcontract.

De hoofdvraag krijgt een voorlopig ontwerpantwoord met 23 requirements, negen kandidaatprincipes, vijftien analyseviews en een partieel geformaliseerd contextmodel. Deze aantallen zijn geen maat voor wetenschappelijke bijdrage. De verbindingsgraaf, mappingcontext en versiegebonden productmanifesten vormen de voorgestelde integratie. De traceability van deelvragen naar toekomstige metrieken maakt die integratie controleerbaar als onderzoeksvoorstel.

De paper stelt niet vast dat de architectuur werkt of dat het bestaande prototype eraan voldoet. Ook voordelen voor menselijke interpretatie, kwaliteitsbeheer en AI blijven hypothesen. Een BAG-fragment begrenst de huidige casusonderbouwing; een gevalideerde keten over meerdere basisregistraties ontbreekt nog. De evaluatie moet vaststellen of de voorgestelde integratie interpretatie en traceerbaarheid verbetert, welke onderdelen daaraan bijdragen en welke extra beheerlast aanvaardbaar is.

\section{Future Work}
De eerstvolgende stap is inhoudelijke bronverdieping: volledige toepasselijke normteksten, een uitgebreidere review van verwante integratiearchitecturen en betrouwbare primaire bronnen voor koppelingen buiten BAG. Daarna worden de formele type- en mappingregels uitgewerkt tot een uitvoerbaar profiel met expliciete uitzonderingen. Een onafhankelijke review moet zowel de juridische context als de ontologische aannames beoordelen.

Vervolgens wordt het prototype feitelijk geïnventariseerd en aan C1--C15 gemapt. De evaluatiecriteria worden vooraf vastgelegd, het testmateriaal wordt onafhankelijk samengesteld en de formatieve episodes worden uitgevoerd. Pas daarna volgt een summatieve vergelijking van bevroren varianten. Een vervolgpublicatie kan die empirische bevindingen rapporteren, inclusief afwezige effecten, onjuist gebleken ontwerpprincipes en noodzakelijke architectuurherzieningen.

Binnen de scope blijven standaardenintegratie, conceptuele modellering, brontraceerbaarheid, governance en een evaluatieplan voor het Stelsel. Buiten de huidige scope vallen een productie-uitrol, migratie van bronregistraties, een volledige juridische beoordeling van concrete verwerkingen, een algemene ontologie voor alle overheidstaken, het trainen van AI-modellen en claims over gemeten maatschappelijke baten. Deze grenzen maken zichtbaar welk vervolgonderzoek nodig is voordat een operationele of generaliseerbare effectclaim kan worden gedaan.
